\section{Introduction}

Periodic products attach multiplicative data to the ordinary cycles
of a polynomial map. For a rational weight $g$, one can ask whether
the product of $g$ is one around every cycle outside a finite
exceptional collection. Even when each individual product is elementary to define,
this is an infinite condition, and the finite set of allowed
exceptions is not supplied as part of the input. A finite test must
therefore address two distinct issues: how to detect all returns
once an exceptional polynomial is fixed, and how to eliminate that
unknown polynomial with a coefficient-independent bound.

This paper resolves both issues for every polynomial of derivative
zero over an algebraic closure of a finite field. The return
polynomials $f^{\circ n}-x$ are then squarefree, even when the
ordinary period is divisible by the characteristic. This turns
polynomial annihilation into an exact statement about ordinary
points. A finite-state representation controls the return tests for
a fixed annihilator. Under the cofinite product condition $\CP$,
a separate rank argument bounds the native period of each visible
exceptional cycle, making a universal
annihilator available. The resulting return bound depends only on
the degree of the map and the height of the rational weight.

Two classical ingredients should be distinguished from this
combination. Coefficient functionals, quotient normal forms, and
residue pairings are established tools; Cattani, Dickenstein, and
Sturmfels give a systematic treatment of normal forms and
multidimensional residues, in particular in Section~4 of the
preprint version of~\cite{CDS1996}. We do not import a complex-analytic residue
formula into positive characteristic: the locally finite Laurent
extractor needed here is proved directly. Finite word cutoffs for
finite-dimensional weighted representations are also classical.
The arbitrary-field formulation in Section~3 and Proposition~3.1
of Kiefer et al.~\cite{Kiefer2013} supplies the relevant context;
we include the elementary matrix-span proof rather than use a
rational-field Gram argument or a complexity theorem.

The specific work is the general-polynomial coefficient and carry
construction for two whole rational-weight products, followed by
the elimination of an unknown finite exceptional set. The split
rank calculation uses two full-rank evaluation maps. Persistence
uses finite multiplicative orders only to keep a specified visible
cycle visible at arbitrarily long returns; it does not require
simultaneous visibility of every exceptional point. These details
are essential to the degree-only conclusion.

Section~\ref{sec:statement} states the finite criterion and handles
actual conjugacy for nonmonic maps. Section~\ref{sec:cycles}
identifies the ordinary-cycle condition with an annihilator ideal
and separates product visibility from derivative filtering.
Sections~\ref{sec:extractor} and~\ref{sec:transfer} prove the
coefficient representation and fixed-annihilator theorem.
Section~\ref{sec:rank} bounds visible native periods, and
Section~\ref{sec:decision} completes the finite decision.
